% !TeX root = ../main.tex
% =============================================================================
% 第 2 章 涉及的理论与技术基础（修订版）
% 修订要点：
% 1. 全章中文全角标点统一
% 2. 新增 §2.4 联邦学习与聚类联邦学习
% 3. 关键公式补充 \label{}，便于跨章引用
% 4. §2.5（原 §2.4）小节号顺延
% =============================================================================
\chapter{相关理论与技术基础}
\label{chap:ch2_background}

本章系统梳理本研究所涉及的天基网络计算卸载、区块链可信协同、联邦学习以及多智能体强化学习四类核心技术的基本原理、典型方法与应用边界，旨在为后续章节中天基协同计算系统的统一建模、长时间尺度协同优化以及短时间尺度卸载决策提供完整的理论支撑与技术依据。天基协同计算系统具有节点高速运动、链路时变、资源异构与可信环境受限等鲜明特征，传统地面边缘计算中所沿用的卸载方法与协同方法难以同时兼顾通信约束的精准刻画与协同决策的可信执行。本章梳理的四类前沿技术，正是本研究实现天基场景下高效卸载与可信协同的核心技术基础。

% -----------------------------------------------------------------------------
\section{天基网络与计算卸载技术}
\label{sec:ch2_satellite_offloading}

天基网络是指以低轨（Low Earth Orbit，LEO）、中轨（Medium Earth Orbit，MEO）和高轨（Geostationary Earth Orbit，GEO）卫星星座为核心，辅以地面控制中心、信关站及用户终端构成的一类天地一体化通信与计算网络。近年来，随着 LEO 巨型星座的快速部署与星上计算能力的持续提升，天基网络已从单纯的“通信中继”角色逐步演化为面向 6G 的“通信—感知—计算”一体化基础设施\cite{He2021,Wu2024ISL,Wang2024NTN,Saad2024NTNRel1718,Mahboob2024AINTN,Wang2025NTN6G}。本节将依次梳理天基网络的层级架构与链路特性、计算卸载的基本范式以及面向 LEO 环境的卸载决策方法。

\subsection{天基网络的层级架构与链路特性}
\label{subsec:ch2_sagin_arch}

目前关于天基网络的体系结构尚无统一规范，学术界普遍采用空—天—地一体化网络（Space–Air–Ground Integrated Network，SAGIN）的层级框架对其进行描述\cite{Wu2024ISL,LMCECN2025,Huang2024ICNLSMC}。在该框架下，LEO 卫星因其轨道低、传播时延小、链路增益高的优势，成为承载星上边缘计算的主要平台\cite{Sun2024MCO,Zhou2024MobilityAware}。本文将天基协同计算网络抽象为如下三层结构：

具体而言，用户终端层主要负责试验任务接入、任务属性提取和结果接收；星上边缘层由具备计算、缓存和星间通信能力的卫星边缘节点构成，承担任务本地处理、协同卸载和资源共享功能；地面控制层由地面控制中心和信关站等组成，主要负责系统管理、模型初始化、策略评估和长期运行维护。上述三层结构与第\ref{chap:ch3_unified_model}章中的 ST、SEN 和 CC 三类节点角色相对应。

天基网络的链路特性显著区别于地面网络，主要体现在三方面。第一，链路时变性强。LEO 卫星轨道周期约为90到100分钟，星—地、星—星之间的可见性窗口在分钟级尺度反复出现与消失，导致链路拓扑随时间剧烈变化\cite{Sun2024MCO}。第二，带宽与时延非对称。星—地下行链路普遍优于上行链路，且馈电链路与用户链路在频段、功率预算上差异显著\cite{Wang2024NTN}。第三，可达性受星历强约束，任意一对节点之间的端到端时延并非常数，而是随轨道相位的周期性变化的随机过程。上述特性使得传统假设链路恒定或缓变的卸载方法难以直接迁移，必须显式建模链路的时变性与可见性窗口约束。

本节梳理的天基网络层级结构是第\ref{chap:ch3_unified_model}章建模工作的应用背景。第\ref{chap:ch3_unified_model}章中将“用户终端层—星上边缘层—地面控制层”三层结构进一步形式化为卫星终端（Satellite Terminal，ST）、卫星边缘节点（Satellite Edge Node，SEN）和地面控制中心（Control Center，CC）三类节点角色，并在此基础上建立链路时延、可达性与资源异构的数学模型。

\subsection{计算卸载基本范式}
\label{subsec:ch2_offloading_paradigm}

计算卸载是指将计算密集型或时延敏感型任务从资源受限的终端节点迁移至算力充裕的边缘或云端节点执行的一类技术，其核心目标是在用户体验质量与系统资源开销之间取得平衡\cite{ref2,Tang2022DRL,Adu2026MADRL}。按卸载粒度划分，卸载方案通常包括二元卸载、部分卸载与依赖型卸载三类。二元卸载将任务视为不可分割的整体，决策变量仅为本地执行或卸载执行；部分卸载允许将任务按比例分割后并行处理；依赖型卸载则进一步考虑任务内部子模块之间的依赖关系，通常以有向无环图（Directed Acyclic Graph，DAG）的形式描述\cite{Guo2024E2E,Qu2024E2E}。

无论采用何种卸载粒度，卸载决策的形式化通常围绕时延、能耗与可靠性三项指标展开。任务在节点 $i$ 上的本地执行时延 $T^{\mathrm{loc}}_i$ 可表示为：
\begin{equation}
  T^{\mathrm{loc}}_i = \frac{c_k}{f_i}
  \label{eq:ch2_local_delay}
\end{equation}
其中 $c_k$ 为任务所需的 CPU 周期数，$f_i$ 为节点 $i$ 的计算频率。当任务被卸载至边缘节点 $j$ 时，端到端时延由传输时延、排队时延与边缘执行时延共同构成：
\begin{equation}
  T^{\mathrm{off}}_{i,j} = \frac{d_k}{r_{i,j}} + T^{\mathrm{queue}}_j + \frac{c_k}{f_j}
  \label{eq:ch2_offload_delay}
\end{equation}
其中 $d_k$ 为任务输入数据量，$r_{i,j}$ 为节点 $i$ 到节点 $j$ 的瞬时传输速率，$T^{\mathrm{queue}}_j$ 为目标节点的排队时延。能耗模型通常采用动态电压频率调整（Dynamic Voltage and Frequency Scaling，DVFS）假设，以芯片功耗系数 $\kappa$ 与计算频率立方的乘积形式表达\cite{Sun2024MCO}。

针对天基场景，计算卸载的研究焦点逐步从地面单跳卸载演进为多跳协同卸载与星—地联合卸载。Sun 等人\cite{Sun2024MCO}考虑 LEO 卫星运动性带来的链路时变，提出了运动感知的卸载策略，联合最小化网络时延与能耗。Wang 等人\cite{LMCECN2025}进一步将主从星协同纳入决策空间，提出多目标 D3QN 算法以兼顾任务完成率、负载均衡、时延与能耗。Guo 等人\cite{Guo2024E2E}则针对端到端任务，将卸载与路由联合建模为 NP 难问题并设计分解算法。近期面向对等卸载、服务部署、地星协同联邦学习和卫星边缘推理的研究进一步拓展了卸载决策的建模边界\cite{Zhang2024PeerOffloading,Tang2024TwoTimescaleJSAC,Zhai2024SECO,ref26,Fan2025SEI,Feng2026LatencyAware}。这些工作虽然在仿真层面验证了卸载在天基场景下的可行性，但普遍假设链路状态、节点信誉等关键变量集中可得，尚未充分考虑大规模、多管理域场景下“信息可信共享”这一前提。

本研究在第\ref{chap:ch3_unified_model}章定义的任务模型基础上，在第\ref{chap:ch5_task_offloading}章构建基于 DQN 的任务层卸载决策方法，并在此基础上引入由第\ref{chap:ch4_system_resource_allocation}章系统层下发的候选约束，以缓解上述工作中“集中可得信息”假设的局限。

\subsection{深度强化学习在卸载决策中的应用}
\label{subsec:ch2_drl_offloading}

由于卸载决策本质上是序列化决策问题，且环境状态具有高维、时变与部分可观特征，基于深度强化学习（Deep Reinforcement Learning，DRL）的方法近年来成为该领域的主流研究范式\cite{Mnih2015dqn,VanHasselt2016DoubleDQN,Wang2016Dueling,Tang2022DRL,Lim2022DRLOS}。DRL 通过将神经网络作为价值函数或策略函数的近似器，结合马尔可夫决策过程（Markov Decision Process，MDP）的形式化框架，实现对卸载策略的端到端学习。

深度 Q 网络（Deep Q-Network，DQN）是该领域的奠基性工作，其核心思想是用神经网络 $Q(s, a; \theta)$ 近似最优动作—价值函数，并通过经验回放和目标网络两项机制缓解训练发散问题\cite{Mnih2015dqn}。DQN 的更新目标为：
\begin{equation}
  y = r + \gamma \max_{a'} Q(s', a'; \theta^-)
  \label{eq:ch2_dqn_target}
\end{equation}
其中 $\theta^-$ 表示目标网络参数。然而原始 DQN 存在过估计偏差与价值估计与动作选择耦合两个问题。Double DQN 通过将动作选择与价值评估解耦缓解了过估计偏差，Dueling DQN 进一步将 Q 函数分解为状态价值与动作优势两支路，提升了在动作差异较小场景下的样本效率\cite{VanHasselt2016DoubleDQN,Wang2016Dueling,Lim2022DRLOS}：
\begin{equation}
  Q(s, a) = V(s) + A(s, a) - \frac{1}{|\mathcal{A}|}\sum_{a'} A(s, a')
  \label{eq:ch2_dueling}
\end{equation}

将 DQN 类方法应用于计算卸载场景的关键在于状态空间与动作空间的合理设计。Tang 与 Wong\cite{Tang2022DRL} 将节点队列状态、信道质量与历史决策共同纳入状态空间，采用 LSTM 与 Dueling DQN 结合的方式建模长期回报。Adu 等人\cite{Adu2026MADRL} 进一步在多接入边缘计算场景下将 DQN 扩展为多智能体形式，使每个边缘设备可基于本地观测做出去中心化卸载决策。Wang 等人\cite{LMCECN2025} 则针对 LEO 巨型星座设计了 D3QN 算法，在多目标（完成率、负载、时延、能耗）联合优化中取得了显著效果。

DQN 类方法是本研究第\ref{chap:ch5_task_offloading}章任务层方法的核心技术路径。第\ref{chap:ch5_task_offloading}章将以轻量化 DQN 为基础，以第\ref{chap:ch3_unified_model}章定义的链上摘要、链路状态与候选约束为状态输入，输出短时间尺度的离散卸载动作，并通过候选约束机制将系统层的协同意图传递到任务层执行。

% -----------------------------------------------------------------------------
\section{区块链与可信协同技术}
\label{sec:ch2_blockchain}

区块链是一种以链式数据结构存储交易、以分布式共识维护账本一致性、以密码学保证数据不可篡改的分布式账本技术\cite{ref3,ref13,ETSI2024BCFL,Villegas2025BETAC}。在天基协同计算场景下，引入区块链可在缺乏可信中心的条件下，为多管理域之间的状态共享、行为审计与模型治理提供透明、可追溯的可信底座\cite{ref12,ref15,ZhangLEO2024,Elmahallawy2025}。本节将分别介绍联盟链与共识机制、节点信誉机制，以及区块链辅助下的多方协同范式。

\subsection{联盟链与共识机制}
\label{subsec:ch2_consortium}

按节点准入策略划分，区块链可分为公有链、联盟链与私有链三类。公有链节点完全开放、采用工作量证明（Proof of Work，PoW）等高资源消耗共识，适合无中心化加密货币场景但难以适配资源受限的天基环境；私有链节点完全封闭、性能高但缺乏多方信任价值；联盟链则在两者之间取得平衡，通过组织级准入机制控制节点身份，采用拜占庭类共识保障一致性，既保留了分布式信任的优势，又显著降低了共识能耗，因而成为面向行业级可信协同的主流选择\cite{ETSI2024BCFL,NetTopoBFT2025}。

实用拜占庭容错（Practical Byzantine Fault Tolerance，PBFT）及其衍生算法是联盟链中最为广泛使用的共识机制。PBFT 通过预准备、准备、提交三阶段的多轮消息交互，在不超过 $\lfloor (n-1)/3 \rfloor$ 个拜占庭节点的前提下保证账本一致性，具有确定性最终性与无需挖矿的优势\cite{Lamport1982Byzantine,Castro1999,NetTopoBFT2025}。然而，经典 PBFT 算法的通信复杂度为
\begin{equation}
  \mathcal{C}_{\mathrm{PBFT}} = \mathcal{O}(n^2),
  \label{eq:ch2_pbft_complexity}
\end{equation}
在节点规模较大时延迟随节点数迅速恶化。为此，后续研究从两个方向加以改进：一是引入分片（Sharding）机制将共识负载分散到多个子链，显著提升系统吞吐\cite{ZhangLEO2024}；二是结合节点信誉对参与共识的节点子集进行动态筛选，将式~（\ref{eq:ch2_pbft_complexity}） 压缩为面向高信誉节点的 $\mathcal{O}(k^2)$ 通信（$k \ll n$），从而在大规模异构网络中保持可用的共识时延\cite{NetTopoBFT2025}。

智能合约（Smart Contract）是联盟链的核心可编程接口，其本质是部署在区块链上、由共识节点协同执行、执行结果上链可追溯的一段确定性代码。智能合约通过将业务逻辑显式表达为代码，使得多方协同中的关键判定（如准入审核、积分发放、违规惩戒）能够在无需可信第三方的条件下自动完成，极大提升了多方协同的执行效率与可审计性\cite{Androulaki2018Fabric,Villegas2025BETAC}。

本研究在第\ref{chap:ch3_unified_model}章中以联盟链作为天基协同计算系统的可信底座，将链上摘要、生命周期事件、节点信誉与模型版本等关键状态以分层共识方式维护；PBFT 类共识与智能合约则在第\ref{chap:ch4_system_resource_allocation}、\ref{chap:ch5_task_offloading}章方法中作为约束执行与凭据流转的实现假设。

\subsection{可信状态共享与节点信誉}
\label{subsec:ch2_reputation}

在多方协同的可信场景下，仅依靠共识机制保证账本一致性尚不足以支撑长期协作，还需进一步引入节点信誉（Reputation）机制对参与方的历史行为进行定量评估，以激励诚实行为、抑制恶意或低质量行为\cite{NetTopoBFT2025,ETSI2024BCFL}。节点信誉的核心在于将节点过往的行为（如任务执行准时率、模型更新质量、共识参与稳定性等）转换为一个随时间演化的可比较标量，并将该标量作为节点参与协同的权重或资格的判据。

典型的信誉演化模型采用指数加权移动平均的形式：
\begin{equation}
  R_i(t+1) = (1-\eta) R_i(t) + \eta r_i(t)
  \label{eq:ch2_reputation_update}
\end{equation}
其中 $R_i(t)$ 为节点 $i$ 在时刻 $t$ 的信誉值，$r_i(t)$ 为本轮的行为评分，$\eta$ 为学习率。信誉机制可与共识算法深度耦合，例如基于信誉的 PBFT 仅允许信誉高于阈值的节点参与共识投票，从而在保证安全的前提下进一步压缩共识规模\cite{NetTopoBFT2025}。

节点信誉机制是第\ref{chap:ch3_unified_model}章中“链上摘要—行为审计—权限调整”反馈环的关键中间变量，亦是第\ref{chap:ch4_system_resource_allocation}章系统层方法在筛选候选节点、第\ref{chap:ch5_task_offloading}章任务层方法在过滤目标节点时所共享的状态接口。式~（\ref{eq:ch2_reputation_update}） 是第\ref{chap:ch3_unified_model}章信誉更新规则的理论原型。

\subsection{区块链辅助的多方协同范式}
\label{subsec:ch2_bc_assisted}

区块链作为可信中介，已被广泛应用于多方协同场景下的去中心化数据共享、可追溯任务调度与多管理域治理。其典型应用范式包括：（1）将关键事件以哈希摘要的形式上链，原始数据保存于链下，形成“链上摘要 + 链下实体”的轻量化存证机制；（2）利用智能合约实现可编程的多方业务规则，例如自动化的资源分配、信誉惩戒和异常处置；（3）将分布式共识与多方协同决策深度耦合，使协同结果具备可审计性\cite{Feng2022,Moscato2023,ETSI2024BCFL,Elmahallawy2025}。

针对天基场景，Zhang 等人\cite{ZhangLEO2024} 提出了面向 LEO 网络的分片化区块链框架，通过将卫星节点按轨道平面划分为多个分片以缓解全网共识开销；Elmahallawy 等人\cite{Elmahallawy2025} 进一步将其拓展至多厂商星座场景，引入区块链辅助机制以应对跨星座间的信任薄弱问题。类似地，面向 LEO 卫星网络的分片化安全联邦学习框架也说明了共识负载拆分与模型更新存证之间的耦合价值\cite{ref29}。这些工作共同表明，区块链不仅是可信存储工具，还能与上层协同决策深度集成，构成多管理域天基协同的基础设施。

本研究将上述范式具体化为面向天基协同计算的“可信控制平面”：第\ref{chap:ch3_unified_model}章定义链上数据集合 $\mathcal{D}_{\mathrm{on}}$ 与摘要—对象哈希映射；第\ref{chap:ch4_system_resource_allocation}章利用链上公共摘要支撑系统层多智能体协同；第\ref{chap:ch5_task_offloading}章利用链上模型版本凭据支撑任务层并发开销预测模型的可信调用。

% -----------------------------------------------------------------------------
\section{联邦学习与聚类联邦学习}
\label{sec:ch2_federated_learning}

随着天基网络中的节点数量、任务多样性与隐私敏感性持续上升，传统的集中式机器学习范式难以满足“数据不出域”与“模型可协同”的双重要求。联邦学习（Federated Learning，FL）作为一类典型的隐私友好型分布式学习框架，能够在保留原始数据于本地的前提下，通过模型参数的多轮交换实现跨节点协同建模，已经成为天基边缘智能与跨星座协同的重要技术基础\cite{McMahan2017,Bonawitz2019,LiFedAvg2020,Li2020FederatedChallenges,Kairouz2021Advances}。本节首先介绍 FedAvg 算法及其在异质性条件下的局限，随后讨论聚类联邦学习的基本思想，最后回到区块链辅助的联邦学习范式与本研究的关联。

\subsection{FedAvg 算法与基本流程}
\label{subsec:ch2_fedavg}

FedAvg 是联邦学习领域最具代表性的基线算法。设全网共有 $N$ 个客户端，第 $k$ 个客户端的本地数据集为 $\mathcal{D}_k$，本地经验风险为 $F_k(w)$，则 FedAvg 的全局优化目标可写为
\begin{equation}
  \min_{w} f(w) = \sum_{k=1}^{N} p_k F_k(w),\qquad
  p_k = \frac{|\mathcal{D}_k|}{\sum_{j=1}^{N} |\mathcal{D}_j|},
  \label{eq:ch2_fedavg_objective}
\end{equation}
其中 $p_k$ 为客户端权重，通常正比于其本地样本量。FedAvg 的工作流程可概括为四步：（1）服务器初始化全局模型 $w^{(0)}$ 并下发至客户端；（2）每个客户端在本地以多轮 SGD 更新得到本地参数 $w_k^{(t+1)}$；（3）客户端将本地参数上传至服务器；（4）服务器按式~（\ref{eq:ch2_fedavg_objective}） 中的权重对客户端参数进行加权平均，得到新的全局模型
\begin{equation}
  w^{(t+1)} = \sum_{k=1}^{N} p_k\, w_k^{(t+1)}.
  \label{eq:ch2_fedavg_aggregation}
\end{equation}
在通信—聚合循环若干轮后，全局模型逐步收敛。FedAvg 的优势在于无需上传原始数据，仅交换模型参数，因而通信开销与隐私风险均较集中式训练显著下降\cite{McMahan2017,LiFedAvg2020}。

\subsection{数据异质性挑战与聚类联邦学习}
\label{subsec:ch2_clustered_fl}

然而，FedAvg 在面对客户端数据非独立同分布（Non-IID）时存在收敛慢、模型质量下降甚至训练发散的问题\cite{Li2020FedProx,Karimireddy2020Scaffold,Zou2024FL,FedFusion2024}。其根本原因在于：当不同客户端的本地最优解 $w_k^*$ 相差较大时，单一全局模型 $w$ 必然会落在某种“折中”位置，对多数客户端而言并非最优。在天基场景下，这一问题尤为突出——不同卫星节点的硬件能力、负载模式、任务类型与可见窗口差异显著，本地数据分布天然具有强异质性。

为缓解这一问题，聚类联邦学习（Clustered Federated Learning，Clustered FL）应运而生。其核心思想是放弃“全网共享单一模型”的强假设，转而将客户端按数据分布或参数更新方向划分为若干簇，每簇维护一个独立的簇模型，仅在簇内进行参数聚合。形式化地，设客户端被划分为 $K$ 个互不相交的簇 $\{\mathcal{C}_1, \mathcal{C}_2, \ldots, \mathcal{C}_K\}$，则第 $c$ 个簇在第 $t+1$ 轮的模型聚合可写为
\begin{equation}
  w_c^{(t+1)} = \sum_{k \in \mathcal{C}_c} \frac{|\mathcal{D}_k|}{\sum_{j \in \mathcal{C}_c} |\mathcal{D}_j|}\, w_k^{(t+1)}.
  \label{eq:ch2_clustered_fedavg}
\end{equation}
与式~（\ref{eq:ch2_fedavg_aggregation}） 相比，式~（\ref{eq:ch2_clustered_fedavg}） 的关键差异在于聚合范围被限制在“数据分布相近”的客户端子集内，从而能够保留局部模式、避免被全局平均效应淹没。簇划分本身可以基于客户端的负载特征、参数更新方向、梯度相似度或硬件能力等多种依据，常用方法包括 $k$-means 聚类、层次聚类与基于参数余弦相似度的动态分组\cite{sattler2021clusteredfl,Zou2024FL,FedFusion2024}。聚类联邦学习已在医疗影像、移动终端与物联网等多个异质性强的场景中表现出优于 FedAvg 的预测精度与收敛稳定性。

\subsection{区块链辅助的联邦学习}
\label{subsec:ch2_bcfl}

虽然联邦学习能够避免原始数据的集中传输，但其安全性仍依赖于一个可信的中心聚合服务器：一旦该服务器失效或被攻陷，恶意更新可能在不被察觉的情况下污染全局模型。区块链辅助的联邦学习（Blockchain-assisted Federated Learning，BCFL）将区块链作为可信聚合中介或可信凭据账本，将每轮的模型更新、参与方贡献、聚合凭据等以哈希摘要的形式上链，既消除了对中心聚合服务器的单点信任依赖，又借助区块链的可追溯性使恶意更新可被识别与回溯\cite{Feng2022,ETSI2024BCFL,Elmahallawy2025}。针对天基场景，Zhang 等人\cite{ZhangLEO2024}提出了面向 LEO 网络的分片化 BCFL 框架，通过将卫星节点按轨道平面划分为多个分片以缓解全网共识开销；Elmahallawy 等人\cite{Elmahallawy2025}进一步将其拓展至多厂商星座场景，引入 BCFL 以应对跨星座间的信任薄弱问题。面向地星融合网络的协同联邦学习研究也表明，局部计算、数据卸载与模型聚合需要联合设计\cite{ref26}。

本节梳理的 FedAvg、聚类联邦学习与 BCFL 是第\ref{chap:ch5_task_offloading}章任务层方法的核心技术路径之一。具体而言，第\ref{chap:ch5_task_offloading}章将基于式~（\ref{eq:ch2_clustered_fedavg}） 的簇内聚合规则训练并发开销预测模型，按节点并发任务数对节点进行分簇以应对工作负载的异质性；同时，将每轮聚合后的模型版本摘要按 BCFL 思想写入第\ref{chap:ch3_unified_model}章定义的模型版本集合 $\mathcal{D}^{\mathrm{model}}_{\mathrm{on}}$，实现“训练—发布—回滚”的可信治理闭环。

% -----------------------------------------------------------------------------
\section{多智能体强化学习与协同决策}
\label{sec:ch2_marl}

天基协同计算系统中，多个 SEN 节点需要在通信受限、信息部分可观的环境下做出长时间尺度的资源分配与候选筛选决策，其本质属于多智能体协同决策问题。与单智能体方法相比，多智能体强化学习（Multi-Agent Reinforcement Learning，MARL）需要同时应对环境非平稳性与动作空间指数级膨胀两大挑战\cite{Shapley1953,Lowe2017,Amato2024CTDE,Liu2025MARL}。本节将依次介绍 MARL 的基本框架、CTDE 训练范式以及代表性算法。

\subsection{多智能体马尔可夫决策过程}
\label{subsec:ch2_dec_pomdp}

多智能体协同决策问题通常形式化为去中心化部分可观马尔可夫决策过程（Decentralized Partially Observable Markov Decision Process，Dec-POMDP），其形式化定义为如下七元组：
\begin{equation}
  \langle \mathcal{N}, \mathcal{S}, \{\mathcal{A}_i\}_{i \in \mathcal{N}}, \mathcal{P},
  \mathcal{R}, \{\mathcal{O}_i\}_{i \in \mathcal{N}}, \gamma \rangle
  \label{eq:ch2_dec_pomdp}
\end{equation}
其中 $\mathcal{N}$ 为智能体集合，$\mathcal{S}$ 为全局状态空间，$\mathcal{A}_i$ 为智能体 $i$ 的动作空间，$\mathcal{P}$ 为状态转移函数，$\mathcal{R}$ 为团队共享奖励函数，$\mathcal{O}_i$ 为智能体 $i$ 的局部观测空间，$\gamma$ 为折扣因子。每个智能体 $i$ 仅可观测其局部观测 $o_i \in \mathcal{O}_i$，根据策略 $\pi_i(a_i \mid o_i)$ 选择动作，所有智能体共同最大化期望累积折扣奖励：
\begin{equation}
  J(\boldsymbol{\pi}) = \mathbb{E}_{\boldsymbol{\pi}}
  \!\left[\sum_{t=0}^{\infty} \gamma^t \mathcal{R}(s_t, \mathbf{a}_t)\right]
  \label{eq:ch2_marl_objective}
\end{equation}

Dec-POMDP 的核心难点在于，从任意单个智能体的角度看，其他智能体策略的演变会使得环境呈现非平稳特征（Non-stationarity），从而违反单智能体强化学习中环境动态平稳的关键假设\cite{Tsitsiklis1997,Liu2025MARL}。这种非平稳性是直接将单智能体方法独立部署于多智能体环境时，常常出现训练不稳定甚至发散的根本原因。

\subsection{集中训练分散执行范式}
\label{subsec:ch2_ctde}

集中训练分散执行（Centralized Training with Decentralized Execution，CTDE）是当前缓解 MARL 非平稳性问题的主流框架\cite{Amato2024CTDE}。CTDE 的核心思想是：在训练阶段，各智能体可访问全局状态或其他智能体的私有信息，从而在更准确的环境模型下学习策略；而在执行阶段，各智能体仅依据本地观测做出决策，从而满足现实通信受限场景的部署约束。CTDE 框架与全局集中式方法相比，显著缓解了动作空间随智能体数指数增长的扩展性问题；与全分散式方法相比，又通过训练阶段的全局信息共享避免了非平稳性带来的训练失稳。

CTDE 范式下的代表性算法可分为价值分解（Value Decomposition）与集中评论家（Centralized Critic）两大类。价值分解方法将团队联合 $Q$ 函数分解为各智能体局部 $Q$ 函数的某种聚合，代表性工作包括 VDN、QMIX 与 QPLEX，其单调性约束保证了局部最优与全局最优的一致性\cite{Sunehag2018VDN,Rashid2018QMIX,Amato2024CTDE}。集中评论家方法则在训练时维护一个以全局状态为输入的中心化评论家网络，并以此指导各智能体局部策略的更新，代表性工作包括 MADDPG 与 MAPPO\cite{Lowe2017,ref5,Liu2025MARL}。

MAPPO 作为近年来最受关注的 CTDE 算法之一，将单智能体 PPO 扩展到多智能体场景，通过共享参数与集中评论家显著提升训练稳定性。其策略更新目标为：
\begin{equation}
  L^{\mathrm{CLIP}}(\theta) = \mathbb{E}\!\left[
    \min\!\Big(\rho_t(\theta) \hat{A}_t,
    \mathrm{clip}(\rho_t(\theta), 1-\epsilon, 1+\epsilon) \hat{A}_t\Big)
    \right]
  \label{eq:ch2_ppo_clip}
\end{equation}
其中 $\rho_t(\theta) = \pi_\theta(a_t \mid o_t) / \pi_{\theta_{\mathrm{old}}}(a_t \mid o_t)$ 为重要性比率，$\hat{A}_t$ 为基于集中评论家估计的优势函数，$\epsilon$ 为裁剪系数。MAPPO 在 SMAC、GRF 等基准上的表现已被多项研究证实其综合性能优于 QMIX 等价值分解方法\cite{ref5}。

\subsection{协同决策中的奖励设计与扩展性}
\label{subsec:ch2_marl_reward}

奖励函数设计是 MARL 应用于具体协同问题的关键环节。共享团队奖励能够在原则上鼓励合作，但在大规模智能体场景下会引发信用分配（Credit Assignment）问题——即难以从团队整体奖励中区分各个智能体的具体贡献。COMA 等反事实基线方法通过构造“假设智能体 $i$ 执行其他动作”的对照基线，显式估计每个智能体对团队奖励的边际贡献，从而缓解信用分配难题\cite{ref11,Amato2024CTDE}。在工程实践中，常用方法是将团队奖励与个体奖励按比例加权，从而在协同性与个体可控性之间寻求平衡。

针对扩展性问题，近期研究探索了基于图神经网络的层级注意力机制\cite{Li2024MARL_GA}，通过对智能体之间的依赖关系建模，将注意力计算压缩到必要的邻居子集，从而显著降低智能体规模带来的状态拼接维度。该思路与本研究中“基于信誉与可见性筛选候选邻居”的协同模式具有较强的契合性。

CTDE 范式与 MAPPO 算法是本研究第\ref{chap:ch4_system_resource_allocation}章系统层方法的核心技术路径。第\ref{chap:ch4_system_resource_allocation}章将以每个 SEN 抽象为一个智能体，以第\ref{chap:ch3_unified_model}章定义的链上共享状态作为观测输入，通过 MAPPO 联合学习任务流向偏好、资源分配比例和外部任务接收意愿等长时间尺度协同动作，并通过候选约束机制将协同意图下发至第\ref{chap:ch5_task_offloading}章的任务层执行。式~（\ref{eq:ch2_ppo_clip}） 即为第\ref{chap:ch4_system_resource_allocation}章策略更新目标的理论原型。

% -----------------------------------------------------------------------------
% -----------------------------------------------------------------------------
\section{本章小结}
\label{sec:ch2_summary}

本章围绕本文后续方法设计所需的理论基础，分别介绍了天基网络与计算卸载、区块链可信协同、联邦学习与聚类联邦学习、多智能体强化学习等相关技术。在天基网络与计算卸载方面，本章梳理了天基网络的层级结构、链路时变特性以及卸载决策的基本时延模型，为第\ref{chap:ch3_unified_model}章构建卫星终端、卫星边缘节点、地面控制中心和动态链路模型提供基础。在区块链可信协同方面，本章介绍了联盟链、共识机制、智能合约和节点信誉机制，为后续可信状态共享、行为审计和链上反馈机制提供支撑。在联邦学习与聚类联邦学习方面，本章说明了分布式数据不出域条件下的模型协同训练方法，为第\ref{chap:ch5_task_offloading}章并发开销预测模型的训练提供理论依据。在多智能体强化学习方面，本章介绍了Dec-POMDP、CTDE和MAPPO等基本方法，为第\ref{chap:ch4_system_resource_allocation}章多节点系统层协同资源分配奠定基础。

总体而言，本章的作用在于为后续章节提供必要的理论和技术支撑，后续第\ref{chap:ch3_unified_model}章将在上述技术基础上建立天基协同计算卸载统一模型；第\ref{chap:ch4_system_resource_allocation}章将结合区块链可信状态共享与多智能体强化学习，研究系统层协同资源分配方法；第\ref{chap:ch5_task_offloading}章将结合聚类联邦学习与DQN，研究并发开销感知的任务层实时卸载方法。由此，本文形成从理论基础到统一建模，再到系统层协同和任务层卸载的递进关系。
